\documentclass[10pt,a4paper]{amsart}
\usepackage[margin=19mm]{geometry}
\usepackage{amsmath,amssymb,mathtools}
\usepackage[scr=rsfs,cal=euler]{mathalfa}
\usepackage{xfrac}
\usepackage[unicode,hidelinks,bookmarks=false]{hyperref}
\newcommand{\ie}{i.e.\ }
\newcommand{\cf}{cf.\ }
\newcommand{\resp}{respectively\ }
\newcommand{\Subsection}{\subsection}
\providecommand{\nobreakdash}{\nobreak\hbox{-}}
\setlength{\emergencystretch}{2em}
\multlinegap0pt
\usepackage{siunitx}
\usepackage{siunitx}
\usepackage{siunitx}
\usepackage{siunitx}
\usepackage{amssymb}
\usepackage{bm}
\usepackage{mathtools}
\usepackage{xcolor}
\usepackage[shortlabels]{enumitem}
\usepackage{float}
\usepackage{algorithm}\usepackage{algpseudocode}
\usepackage{tikz}
\usepackage{csquotes}
\usepackage{siunitx}
\usepackage[normalem]{ulem}
\newtheorem{theorem}{Theorem}
\def\gL{{\mathcal{L}}}
\newcommand{\lrangle}[1]{\langle#1\rangle}
\def\vk{{\bm{k}}}
\def\vq{{\bm{q}}}
\def\vt{{\bm{t}}}
\def\vx{{\bm{x}}}
\title{Kernelized Linear Attention:\\Breaking the Capacity Wall with Symmetric Cones}
\author{Ayoub Ghriss\, , Sourav Chakraborty}
\date{Source: arXiv:2607.17419v1 (CC BY 4.0)}
\begin{document}
\maketitle

\noindent\emph{Source and licence.} Kernelized Linear Attention:\\Breaking the Capacity Wall with Symmetric Cones. Version arXiv:2607.17419v1 (CC BY 4.0), distributed by arXiv under the Creative Commons Attribution 4.0 International licence, http://creativecommons.org/licenses/by/4.0/, as recorded in the arXiv licence metadata for this exact version and shown on the abstract page https://arxiv.org/abs/2607.17419v1. Full licence notice: \texttt{licenses/arxiv-2607.17419-CC-BY-4.0.txt}.\par\smallskip

\noindent\textit{Editorial scope.} Counted unit: the article's theorem thm:softmax-cone-decomp (source0.tex lines 2938--2958). The article's complete original proof of it is delivered with it (source0.tex lines 2960--3017, one \texttt{proof} environment of 58 lines). No mathematics is written, repaired or re-derived: the statement and the proof are the article's own lines, in the article's own notation, and no retained line is substituted or re-flowed.  No further source-local context is needed: every cross-reference of every retained line resolves inside the two delivered spans.  Nothing was omitted from any retained span, so the package carries no omission record.  No retained line contains a citation, so the wrapper carries no bibliography and no bibliography entry.  No retained line contains a figure environment, an image-inclusion command or drawing code, so no image is delivered and the package declares no asset dependency.  Editorial formatting only, in addition: an AMS \texttt{amsart} preamble that restates the article's own declarations and macros used by the retained lines, namely the environments \texttt{theorem} on separate counters, together with the standard packages those lines need.  The retained environments print in this document as Theorem 1 in order of appearance, whereas the article prints the same items with the numbering of its own full document.  The complete native source of the article as distributed in the arXiv e-print for this exact version is bundled unchanged as \texttt{sources/205570/source0.tex}.
\begin{theorem}[Softmax Taylor components factor through cone atoms]
  \label{thm:softmax-cone-decomp}
  Let $z=\lrangle{\vq,\vk}$ for unit vectors $\vq,\vk\in\mathbb{S}^{d-1}$, and define
  $\vt_r(\vx)=\operatorname{vec}(\vx^{\otimes r})$. Every even Taylor monomial factors through
  rank-one PSD atoms:
  \[
    z^{2r}
    =
    \lrangle{\vt_r(\vq)\vt_r(\vq)^\top,
      \vt_r(\vk)\vt_r(\vk)^\top}_F.
  \]
  The odd-degree terms become nonnegative after Lorentz gating, and the exponential kernel admits
  the decomposition:
  \[
    e^z
    =
    \sum_{r=0}^{\infty}\frac{z^{2r}(1+z)}{(2r+1)!}
    +
    \sum_{r=1}^{\infty}\frac{2r\,z^{2r}}{(2r+1)!}.
  \]
\end{theorem}

\begin{proof}

For unit $\vq,\vk$ and $\vt_r(\vx)=\operatorname{vec}(\vx^{\otimes r})$:
\[
  \lrangle{\vt_r(\vq),\vt_r(\vk)}
  = \lrangle{\vq,\vk}^r
  = z^r .
\]
Thus:
\[
  z^{2r}
  =
  \lrangle{\vt_r(\vq),\vt_r(\vk)}^2
  =
  \lrangle{\vt_r(\vq)\vt_r(\vq)^\top,
    \vt_r(\vk)\vt_r(\vk)^\top}_F,
\]
which is the inner product of two rank-one rays in the PSD cone over the $r$-fold tensor feature
space. An odd monomial satisfies $z^{2r+1}<0$ at $\vk=-\vq$ and therefore violates
nonnegativity on the full sphere.

Let $\ell(\vx)=(\vx,1)$. For $\vx\in\mathbb{S}^{d-1}$, $\ell(\vx)\in\partial\gL^{d+1}_{+}$ and:
\[
  \lrangle{\ell(\vq),\ell(\vk)} = 1+z.
\]
Consequently:
\[
  z^{2r}(1+z)
  =
  \lrangle{\vt_r(\vq)\vt_r(\vq)^\top,
    \vt_r(\vk)\vt_r(\vk)^\top}_F
  \lrangle{\ell(\vq),\ell(\vk)},
\]
the product of a rank-one PSD-ray kernel and a Lorentz-ray kernel. Equivalently, define:
\[
  \Phi_r(\vx)
  =
  \left(\vt_r(\vx)\vt_r(\vx)^\top\right)\otimes \ell(\vx).
\]
Then $\lrangle{\Phi_r(\vq),\Phi_r(\vk)}=z^{2r}(1+z)$ under the tensor-product inner product, making
the ``multiplication'' a standard feature-map tensoring operation.

It remains to check that the Taylor coefficients can be regrouped this way. Absolute convergence on
$[-1,1]$ permits rearrangement, and:
\[
  \sum_{r=0}^{\infty}\frac{z^{2r}(1+z)}{(2r+1)!}
  +
  \sum_{r=1}^{\infty}\frac{2r\,z^{2r}}{(2r+1)!}
\]
has odd coefficient $1/(2r+1)!$ on $z^{2r+1}$ and even coefficient:
\[
  \frac{1}{(2r+1)!}+\frac{2r}{(2r+1)!}
  = \frac{1}{(2r)!}
\]
on $z^{2r}$ for $r\ge1$, while the constant term is $1$. This is exactly the Taylor series of
$e^z$.

\end{proof}

\end{document}
